\clearpage
\chapter*{Đính chính cho Báo cáo số~232}
\addcontentsline{toc}{chapter}{Đính chính cho Báo cáo số~232}
\sid{e232}
\markboth{Đính chính 232}{Đính chính 232}
\src{303}

\noindent\textbf{7. TDTE V: Các lược đồ Picard. Các định lý tồn tại.}

\begin{center}
[Séminaire Bourbaki, tập 14, 1961/62, số 232, 19 trang.]
\end{center}

\noindent\textbf{Trang 7, Nhận xét 3.3.} --- Vấn đề nêu ra ở đây đã được
MUMFORD giải đáp khẳng định (xem phần Bình luận cho báo cáo kế tiếp).

\noindent\textbf{Trang 13, Nhận xét 5.2.} --- Như tôi chỉ ra ở đầu báo cáo kế
tiếp, phỏng đoán tồn tại được nêu ở đây là sai; tuy nhiên, bằng các phương pháp
của mình, MUMFORD chứng minh được một định lý yếu hơn đôi chút.

\noindent\textbf{Trang 14, dòng 7.} --- Sau ``đầy đủ'', thêm ``có trường thặng
dư đóng đại số''. Hơn nữa, phản ví dụ của MUMFORD vừa nêu cho thấy hạn chế này
là thiết yếu.

\noindent\textbf{Trang 17, Nhận xét 6.6.} --- Như chúng tôi chỉ ra ở đầu báo
cáo kế tiếp, vấn đề nêu ra ở đây vừa được MURRE giải đáp khẳng định.

\noindent\textbf{Trang 18, Nhận xét 6.7.} --- Vấn đề nêu ra ở đây có lời giải
khẳng định (xem đoạn cuối của phần Bình luận cho báo cáo kế tiếp).

\noindent\textbf{Trang 18, Nhận xét 6.8, dòng thứ 4.} --- Thay ``chính quy''
bằng ``trơn trên $k$''.

\par
